\section{The zero-volume boundary}

The full-dimensional volume identity cannot simply be continued to
multiplier-ideal volume zero.
The next theorem is a separate toric result.
Here \(\Delta(T)\) is the Newton body defined directly from the convex
toric potential.  It agrees with the partial Okounkov body on the
positive-volume locus, but at volume zero it is a separate toric extension
of \Cref{def:ok-algebraic-partial-body}.

\begin{theorem}\label{thm:ok-toric-zero-mass}
Let \(X\) be a smooth projective toric variety.
Let \(T\) be a torus-invariant positive \((1,1)\)-current in an integral or
rational big toric class, and let
\[
                             K=\Delta(T)
\]
be its Newton body, in the separate toric zero-volume sense described
above.  No use of \Cref{def:ok-algebraic-partial-body} at zero volume is
intended.
Then Cao's numerical dimension satisfies
\[
                         \nu_{\mathrm{Cao}}(T)=\dim K.
\tag{55.50}\label{eq:ok-toric-numdim}
\]
\end{theorem}

\begin{proof}
Choose a toric \(\mathbf Q\)-divisor \(D\) representing the class of \(T\).
After multiplication by a common denominator, we may suppose that \(D\) is
integral and that \(T\) is the curvature current of a torus-invariant
semipositive singular metric \(h_0\rme^{-\varphi}\) on
\(\mathcal O_X(D)\).  This operation replaces \(K\) by a positive homothetic
copy and does not change either \(\dim K\) or Cao's numerical dimension.  Let
\(M\) and \(N\) be the character and cocharacter lattices.  We normalize
Lebesgue measure \(\lambda_M\) on \(M_{\mathbf R}\) so that a fundamental
parallelepiped of \(M\) has measure one.

Choose a very ample integral toric divisor \(A\), with full-dimensional
moment polytope \(Q\), and a torus-invariant K\"ahler form \(\omega_A\) in
its class.  We use logarithmic coordinates characterized by
\(\langle u,\operatorname{Log}z\rangle=\log|\chi^u(z)|\), and let
\(F_T\) and \(F_A\)
be the corresponding convex weights of the \emph{norms} on
\(N_{\mathbf R}\).  Thus a local squared-norm weight in the convention of
\Cref{sec:ch51-squared-norm} is twice the corresponding norm weight.  The
Newton body is the closure of the effective domain of the convex conjugate
of \(F_T\).  Since the conjugate of a sum is the
lower-semicontinuous infimal convolution, its closed effective domain is the
Minkowski sum of the two compact effective domains.  Consequently, for
rational \(\varepsilon>0\),
\[
                  \Delta(T+\varepsilon\omega_A)=K+\varepsilon Q.
\tag{55.51}\label{eq:ok-toric-minkowski}
\]

We now prove the section-counting identity needed for Cao's criterion.  Put
\(B_\varepsilon=K+\varepsilon Q\), and take \(m\) sufficiently divisible that
\(m\varepsilon\in\mathbf Z\).  Set
\[
 W_{m,\varepsilon}:=
 H^0\!\left(X,\mathcal O_X(mD+m\varepsilon A)
                 \otimes\mathcal I(m\varphi)\right).
\]
The multiplier ideal is torus invariant, so \(W_{m,\varepsilon}\) is the
direct sum of its one-dimensional character spaces.  We determine their
weights chart by chart.

Let \(\sigma\) be a maximal cone of the smooth fan of \(X\).  Its primitive
ray generators form a basis of \(N\); let
\(e_{\sigma,1},\ldots,e_{\sigma,n}\) be the dual basis of \(M\), and put
\(\kappa_\sigma=\sum_j e_{\sigma,j}\).  Let
\(d_\sigma,a_\sigma\in M\) be the Cartier data of \(D,A\) on the affine chart
\(U_\sigma\).  In the corresponding toric frame, the coefficient of the
character section \(\chi^u\) of
\(\mathcal O_X(mD+m\varepsilon A)\) is the monomial with exponent
\[
       \alpha_\sigma:=u-m d_\sigma-m\varepsilon a_\sigma.
\]
Let \(\Phi_\sigma\) be the local squared-norm psh weight of \(T\) in this
frame, and put \(\psi_\sigma=\Phi_\sigma/2\).  Thus \(\psi_\sigma\) is the
norm weight used by Guenancia.  The global relative potential \(\varphi\)
and \(\Phi_\sigma\) differ by a smooth function, and hence
\[
 \mathcal I(m\varphi)|_{U_\sigma}
   =\mathcal I(m\Phi_\sigma)
   =\mathcal I_{\mathrm G}(m\psi_\sigma),
\]
where the last ideal is defined by local integrability against
\(\rme^{-2m\psi_\sigma}\).  The closure of the local Newton convex set of
\(\psi_\sigma\) is
\[
              \overline{P(\psi_\sigma)}
                    =K-d_\sigma+\sigma^\vee.
\]
Indeed, subtracting the Cartier linear function translates the effective
domain by \(-d_\sigma\), while allowing holomorphic monomials adds the dual
cone.  The set \(P(\psi_\sigma)\) is convex and the displayed closure is
full-dimensional.  A convex set with full-dimensional closure has the same
interior as its closure, and hence
\[
 \operatorname{int}P(m\psi_\sigma)
   =\operatorname{int}\bigl(m(K-d_\sigma)+\sigma^\vee\bigr).
\]

Guenancia's toric multiplier-ideal criterion
\cite[Theorem~1.13]{Guenancia12} says that a holomorphic monomial of exponent
\(\alpha_\sigma\) belongs to \(\mathcal I_{\mathrm G}(m\psi_\sigma)\)
exactly when \(\alpha_\sigma+\kappa_\sigma\) lies in this interior.  This
condition also recovers holomorphicity, so it may be applied to an arbitrary
character \(\chi^u\).  Indeed, if
\[
 P_D:=\bigcap_{\tau\in\Sigma(n)}(d_\tau+\tau^\vee)
\]
is the toric polytope of \(D\), extension of the invariant psh metric across
the toric charts gives \(K\subseteq P_D\).  Hence
\(m(K-d_\sigma)+\sigma^\vee\subseteq\sigma^\vee\).  Membership of
\(\alpha_\sigma+\kappa_\sigma\) in its interior implies, in the smooth chart
basis, \(\alpha_{\sigma,j}+1>0\) for every \(j\).  Since these coordinates
are integers, \(\alpha_{\sigma,j}\geq0\), as required.

After cancelling \(m d_\sigma\), we obtain
\[
 \chi^u\in W_{m,\varepsilon}
 \quad\Longleftrightarrow\quad
 u+\kappa_\sigma\in
 \operatorname{int}(mB_\varepsilon+\sigma^\vee)
 \quad\text{for every maximal }\sigma.
\]
Here we used
\(Q+\sigma^\vee=a_\sigma+\sigma^\vee\), which holds because \(A\) is
very ample.  Notice that the criterion includes both the strict interior and
the fixed anticanonical shift \(\kappa_\sigma\); neither may be discarded at
finite level.

Let \(E_{m,\varepsilon}\subset M\) be the set on the right.  Then
\(\dim W_{m,\varepsilon}=\#E_{m,\varepsilon}\) by the toric character
dictionary \cite[\S4.3]{CLS11}.  Completeness of the fan gives the elementary
gluing identity
\[
 B_\varepsilon
   =\bigcap_{\sigma\in\Sigma(n)}
        (B_\varepsilon+\sigma^\vee).
\]
To see the reverse inclusion, if \(x\notin B_\varepsilon\), strict separation
gives \(v\in N_{\mathbf R}\) such that
\(\langle x,v\rangle<\min_{b\in B_\varepsilon}\langle b,v\rangle\).
Choose a maximal cone containing \(v\).  Every point of
\(B_\varepsilon+\sigma^\vee\) has pairing with \(v\) at least that minimum,
so \(x\) is excluded from that member of the intersection.

After division by \(m\), the exact character conditions cut out
\[
 C_{m,\varepsilon}:=
 \bigcap_{\sigma\in\Sigma(n)}
 \left(-\frac{\kappa_\sigma}{m}
       +\operatorname{int}(B_\varepsilon+\sigma^\vee)\right),
 \qquad
 E_{m,\varepsilon}=mC_{m,\varepsilon}\cap M.
\]
The fan is finite and the shifts are \(O(m^{-1})\).  Hence every compact
subset of \(\operatorname{int}B_\varepsilon\) lies in
\(C_{m,\varepsilon}\) for all large \(m\), while
\(C_{m,\varepsilon}\) lies in any prescribed neighborhood of
\(B_\varepsilon\) for all large \(m\).  For completeness, the second
assertion follows by contradiction: the intersection is uniformly bounded
because cones containing both signs of each coordinate direction give
uniform upper and lower bounds; a convergent sequence staying a fixed
distance from \(B_\varepsilon\) would, after passage to the limit, belong to
every closed set \(B_\varepsilon+\sigma^\vee\), hence to
\(B_\varepsilon\), a contradiction.

Since \(\varepsilon>0\) and \(Q\) has interior, \(B_\varepsilon\) is
full-dimensional and its boundary has \(\lambda_M\)-measure zero.  Choose
\(b\in\operatorname{int}B_\varepsilon\).  The homothetic inner bodies
\(b+(1-\delta)(B_\varepsilon-b)\) and arbitrarily small outer parallel
bodies squeeze \(C_{m,\varepsilon}\).  The lattice Riemann-sum formula for a
bounded convex body \(C\),
\[
             \frac{\#(mC\cap M)}{m^n}\longrightarrow\lambda_M(C),
\]
therefore yields
\[
 \lim_{\substack{m\to\infty\\m\varepsilon\in\mathbf Z}}
 \frac{\dim W_{m,\varepsilon}}{m^n}
       =\lambda_M(K+\varepsilon Q).
\]
One elementary proof of the Riemann-sum formula surrounds each lattice point
by its covolume-\(m^{-n}\) fundamental cell; the symmetric difference is
contained in a shrinking neighborhood of \(\partial C\).  This also explains
why no factor \(n!\) occurs here.  We have proved an actual limit, not a
finite-level equality or merely a limsup.

Since the preceding inner limit exists, Cao's projective criterion
\cite[Proposition~4.4]{Cao14} becomes, in this normalization,
\[
 \lim_{\substack{\varepsilon\downarrow0\\
                  \varepsilon\in\mathbf Q}}
 \frac{\log\!\left(
   \displaystyle\lim_{\substack{m\to\infty\\m\varepsilon\in\mathbf Z}}
   \frac{h^0\!\left(X,\mathcal O_X(mD+m\varepsilon A)
                    \otimes\mathcal I(m\varphi)\right)}{m^n}
   \right)}{\log\varepsilon}
 =n-\nu_{\mathrm{Cao}}(T).
\]
The ordinary multiplier ideal in this formula is exactly the one counted
above; Cao proves that its asymptotics agree with those of the
upper-regularized ideal used during his proof.

Put \(d=\dim K\).
After translation, \(K\) contains a \(d\)-dimensional ball in its affine
span and is contained in a bounded \(d\)-dimensional box.  The body \(Q\)
contains an \(n\)-dimensional ball and is contained in a bounded
\(n\)-dimensional box.  Splitting \(M_{\mathbf R}\) into the direction space
of \(K\) and a complement, and comparing with products of balls and boxes,
gives constants \(c,C>0\) such that
\[
 c\varepsilon^{\,n-d}
 \leq\lambda_M(K+\varepsilon Q)
 \leq C\varepsilon^{\,n-d}
\tag{55.52}\label{eq:ok-toric-tube-growth}
\]
for all sufficiently small \(\varepsilon>0\).
Thus
\(\log\lambda_M(K+\varepsilon Q)
 =(n-d)\log\varepsilon+O(1)\).  Substitution in Cao's criterion gives
\(n-\nu_{\mathrm{Cao}}(T)=n-d\), hence
\(\nu_{\mathrm{Cao}}(T)=d\).
\end{proof}

\begin{remark}[Exact scope]
\label{rem:ok-toric-scope}
\Cref{thm:ok-toric-zero-mass} does not include an irrational real toric
class.
Such an extension needs both a definition of its Newton body and an
approximation theorem compatible with Cao's numerical dimension.
Neither is a formal consequence of \eqref{eq:ok-toric-numdim}.
\end{remark}
